\section{问题一：质量评价与领域配比}
\subsection{质量指标统一与冲突分析}
A1--A3 合计 272505 条记录。对 A1 的 22 项指标按题给语义统一正向与反向，在 A1 的 $1\%$--$99\%$ 分位确定稳健标尺，再原样用于 A2/A3。将指标分为知识价值、语言表达、结构丰富和洁净规范四组。每组先求稳健组分，再用 Huber 位置函数合成相对质量分 $Q_A$；其作用是比较当前附件中的样本，不能解释成单篇文本的真实训练收益。冲突度定义为四组评分六个两两绝对差的平均值；最大的差值对仅用于标注主导冲突方向，以 A1 的第 90 百分位 $0.2714$ 为阈值，随后保持阈值不变。

A1 的领域中位质量分分别为 arxiv 0.782、stackexchange 0.715、commoncrawl 0.707、wikipedia 0.622、github 0.577、book 0.569。c4 的高冲突率为 0.2876；其中“洁净规范高而语言表达低”等方向占主导，驱动指标包括行末标点、专业性和教育价值。该分析指出指标结构的不一致，未将相关性解释为文本语义的因果成因。A2/A3 中 arxiv 和 github 的中位分相对 A1 分别变动 $-0.00532$ 和 $-0.00086$。另从 A1 七域按分位和冲突程度抽取 49 条原文，核对 ID、领域与长度；该样本仅是可追溯的定性审计，不等于独立专家盲评。
\fig{problem1/quality/figures/quality_by_domain.pdf}{0.82}{A1 各领域的相对质量分布。分数仅在当前指标与标尺下可比。}
\fig{problem1/quality/figures/conflict_profile_by_domain.pdf}{0.82}{各领域高冲突比例及结构，阈值由 A1 固定。}

\subsection{17 域配比响应面与模型选择}
配比处于单纯形，直接对各占比做无约束回归会引入和为 1 的共线性。借鉴组成数据分析\cite{aitchison}，将正配比以等距对数比变换 $z=\operatorname{ILR}(p)$ 映射到 16 维，再以 Ridge 稳定估计平均 Loss 响应：
\begin{equation}
\widehat L_{\rm mix}(p)=b_0+b^\top z+z^\top H z,\qquad H=H^\top .
\end{equation}
模型仅在 A4--A5 内部五折交叉验证中选择。二次模型的 MSE 为 0.1471，线性模型为 0.1944；留出的 A6--A7 的 1M 平均 Loss RMSE 为 0.23195，Spearman 为 0.66589。该误差针对每个配方的 13 域平均 Loss；逐域合并 RMSE 为 0.33256，不能混为同一指标。真实 $R^2=1-\mathrm{SSE}/\mathrm{SST}$，独立 1M 测试为 0.3023，不能引用早期误算的 0.946。

A16 的 17 域中只有 6 域可直接或近似对应 A1 的评分域。加入已覆盖域质量特征虽然降低内部 CV 误差，却未同时满足预设的独立测试 2\% 改善门槛，因此最终预测器不加入该质量项。其余 11 域没有独立观测质量分，不以推断代理冒充实测。
\fig{problem1/mixture/figures/mixture_prediction_1m.pdf}{0.83}{1M 独立配比检验：预测与观测的平均 Loss。}

\subsection{推荐配比、非加性与尺度限制}
主推荐取实测标准化 Loss 最优前 10\% 的 52 个配方的质心。占比最高的五域为 pile\_cc 15.13\%、wikipedia\_en 11.15\%、pubmed\_central 10.39\%、stackexchange 8.54\% 和 uspto\_backgrounds 7.99\%。其和为 53.20\%，单域最大份额为 15.13\%。这是下一轮训练的实验候选，不是已验证的全局最优。响应面边界优化解单列为诊断结果，不替代质心。
\fig{problem1/mixture/figures/recommended_mixture.pdf}{0.84}{17 域实测优良配方质心及其占比。}

为讨论组合效应，保存二次 ILR 响应面的系数、Hessian 和按配方行重抽样结果；在训练均值及推荐质心两个锚点，对无序领域对做各增加 0.5 个百分点的联合扰动。136 对无序组合经 BH--FDR 修正后，分别有 30 与 38 对呈模型内非加性。结论依赖锚点、重分配路径和二次响应面，不能称为领域的因果互补。对应的 272 条有向一阶替换仅作局部预测，不赋予独立显著性。

A6--A11 的成对检验表按 1M、60M、1B 分层：60M 与 1B 的平均 Loss 原值 RMSE 分别为 1.5364、2.7544，含有很大的尺度截距偏差；逐尺度去均值后的 RMSE 分别为 0.2160、0.1537，排序相关为 0.6258、0.5266。故 1M 模型移到 60M、1B 后只能部分保持配比效应；在 10B、70B 的估计集上排序相关转负，而这些 Loss 本身又是估计值。因此不用 1M 响应面确定大模型配比，也不将其绝对 Loss 与后续 B6 相加。
